\documentclass[10pt,a4paper]{amsart}
\usepackage[margin=19mm]{geometry}
\usepackage{amsmath,amssymb,mathtools}
\usepackage[scr=rsfs,cal=euler]{mathalfa}
\usepackage{xfrac}
\usepackage[unicode,hidelinks,bookmarks=false]{hyperref}
\newcommand{\ie}{i.e.\ }
\newcommand{\cf}{cf.\ }
\newcommand{\resp}{respectively\ }
\newcommand{\Subsection}{\subsection}
\providecommand{\nobreakdash}{\nobreak\hbox{-}}
\setlength{\emergencystretch}{2em}
\multlinegap0pt
\usepackage{mathtools}
\usepackage{amssymb}
\usepackage{dsfont}
\usepackage{xcolor}
\usepackage[shortlabels]{enumitem}
\newtheorem{theorem}{Theorem}
\newtheorem{remark}{Remark}
\newtheorem{proposition}{Proposition}
\newcommand{\aar}[2][]{\xrightarrow[#1]{#2}}
\title{A Mayer-Vietoris calculus for regular denominators}
\author{Elena Caviglia, Amartya Goswami}
\date{Source: arXiv:2608.21102v1 (CC BY 4.0)}
\begin{document}
\maketitle

\noindent\emph{Source and licence.} A Mayer-Vietoris calculus for regular denominators. Version arXiv:2608.21102v1 (CC BY 4.0), distributed by arXiv under the Creative Commons Attribution 4.0 International licence, http://creativecommons.org/licenses/by/4.0/, as recorded in the arXiv licence metadata for this exact version and shown on the abstract page https://arxiv.org/abs/2608.21102v1. Full licence notice: \texttt{licenses/arxiv-2608.21102-CC-BY-4.0.txt}.\par\smallskip

\noindent\textit{Editorial scope.} Counted unit: the article's proposition propeqfacts (source0.tex lines 197--203). The article's complete original proof of it is delivered with it (source0.tex lines 204--258, one \texttt{proof} environment of 55 lines). No mathematics is written, repaired or re-derived: the statement and the proof are the article's own lines, in the article's own notation, and no retained line is substituted or re-flowed.  Retained uncounted as the source-local context the proof needs: lines 144--161; lines 187--191.  Nothing was omitted from any retained span, so the package carries no omission record.  No retained line contains a citation, so the wrapper carries no bibliography and no bibliography entry.  No retained line contains a figure environment, an image-inclusion command or drawing code, so no image is delivered and the package declares no asset dependency.  Editorial formatting only, in addition: an AMS \texttt{amsart} preamble that restates the article's own declarations and macros used by the retained lines, namely the environments \texttt{theorem}, \texttt{remark}, \texttt{proposition} on separate counters, together with the standard packages those lines need.  The retained environments print in this document as Theorem 1, Remark 1, Proposition 1 in order of appearance, whereas the article prints the same items with the numbering of its own full document.  The complete native source of the article as distributed in the arXiv e-print for this exact version is bundled unchanged as \texttt{sources/208225/source0.tex}.
\begin{theorem}\label{thmMayerVietsequence}
    Let $I,J$ be ideals of $A$ and let $a\in A$. There is an exact sequence 
    $$0\,\aar{} K^{I \cap J}_a \aar{\eta_a} {K_a^I \oplus K_a^J} \aar{\alpha_a} K_a^{I+J}  \aar{\delta_a} C_a^{I \cap J} \aar{\beta_a} {C_a^I \oplus C_a^J} \aar{\gamma_a} C_a^{I+J} \aar{} 0 $$
    where 
\begin{enumerate}
\item[$\bullet$] $\eta_a\bigl(c+(I\cap J)\bigr)
    = (c+I,c+J),$
\item[$\bullet$] $\alpha_a(d+I,b+J)
  =d-b+(I+J),$
\item[$\bullet$]    $\beta_a\bigl(d+(I\cap J+aA)\bigr)
   =(d+(I+aA),d+(J+aA)),$
\item[$\bullet$] $\gamma_a(d+(I+aA),b+(J+aA))
    =d-b+(I+J+aA),$

    \item[$\bullet$] $\delta_a\bigl(c+(I+J)\bigr)=j+(I\cap J+aA)
    $ with $i\in I$ and $j\in J$ such that $ac=i+j$.
\end{enumerate}
\end{theorem}

\begin{remark}\label{remkerdelta}
    By construction, the kernel of the connecting homomorphism $\delta_{a}$ of Theorem \ref{thmMayerVietsequence} is given by 
    $$\ker(\delta_a)=((I:a)+(J:a))/(I+J).$$
    This means that the injectivity of $\delta_a$ is determined by whether $(I:a)+(J:a)=I+J$. This fact will be useful later.
\end{remark}

 \begin{proposition}\label{propeqfacts}
     Let $I,J$ be ideals of $A$ and let $a\in S_{I\cap J}$. The following facts are equivalent:
     \begin{enumerate}
         \item $a\in S_I \cap S_J$;
         \item $(I:a)+(J:a)=I+J$.
     \end{enumerate}
 \end{proposition}

\begin{proof}
	Since $a\in S_{I\cap J}$, we have
	$
	K_a^{I\cap J}=0.
	$
	Hence the exact sequence of Theorem~\ref{thmMayerVietsequence} begins as
	\[
	0\longrightarrow 0
	\longrightarrow K_a^I\oplus K_a^J
	\xrightarrow{\alpha_a} K_a^{I+J}
	\xrightarrow{\delta_a} C_a^{I\cap J}.
	\]
	Exactness at $K_a^I\oplus K_a^J$ shows that $\alpha_a$ is injective, while
	exactness at $K_a^{I+J}$ gives
	$
	\operatorname{Im}(\alpha_a)=\ker(\delta_a).
	$
	Consequently, restriction of the codomain of $\alpha_a$ yields an
	isomorphism
	\[
	K_a^I\oplus K_a^J \cong \ker(\delta_a).
	\]
	
	Now
	$
	a\in S_I\cap S_J
	$
	if and only if
	$
	K_a^I=K_a^J=0,
	$
	which is equivalent to
	$
	K_a^I\oplus K_a^J=0.
	$
	By the preceding isomorphism, this is equivalent to
	$
	\ker(\delta_a)=0.
	$
	By Remark~\ref{remkerdelta},
	\[
	\ker(\delta_a)
	=
	\frac{(I:a)+(J:a)}{I+J}.
	\]
	Therefore,
	$
	\ker(\delta_a)=0
	$
	if and only if
	\[
	(I:a)+(J:a)=I+J.
	\]
	This proves the equivalence.
\end{proof}

\end{document}
